%% ============================================================
%%  07_anhang.tex
%%  Kapitel A: Anhang
%% ============================================================

\chapter{Anhang}
\label{chap:anhang}

% Alles, was den Lesefluss im Hauptteil gestört hätte,
% aber zur Vollständigkeit gehört.

% ----------------------------------------------------------------
\section{Zugang zum Quellcode}
\label{sec:anhang_code}


Der im Rahmen dieser Arbeit entwickelte Quellcode ist öffentlich über GitHub
verfügbar und gliedert sich in zwei Repositories.


Das erste Repository enthält die Implementierung der Abstraktionsschicht
\texttt{qc\_executor}, welche die einheitliche Ansteuerung der verschiedenen
Backends kapselt:


\url{https://github.com/flaqship/qc-executor}


Das zweite Repository umfasst das Benchmarking-System, mit dem die vorgestellten Performance-Messungen
durchgeführt wurden:

\url{https://github.com/dfaigle/benchmarking-system}

% Längere Listings, die im Hauptteil zu viel Platz brauchen.


\newpage
\section{Eidesstattliche Erklärung}
\addcontentsline{toc}{chapter}{Eidesstattliche Erklärung}
\thispagestyle{empty}

Ich erkläre hiermit an Eides statt, dass ich die vorliegende
Bachelor-Thesis mit dem Titel

\begin{center}
  \emph{\ThesisTitel}
\end{center}

selbstständig und ohne Benutzung anderer als der angegebenen Hilfsmittel
angefertigt habe.

Bei der Erstellung dieser Arbeit wurden KI-basierte Werkzeuge unterstützend
eingesetzt, und zwar zur Rechtschreib- und Formulierungshilfe, beim Finden
von Quellen sowie zur Unterstützung beim Programmieren. Die inhaltliche
Verantwortung für die Arbeit liegt vollständig bei mir.

Alle Stellen, die wörtlich oder sinngemäß aus veröffentlichten oder nicht
veröffentlichten Schriften entnommen wurden, sind als solche kenntlich
gemacht. Die Arbeit ist in gleicher oder ähnlicher Form noch nicht als
Prüfungsarbeit eingereicht worden.

Ich bin mir bewusst, dass eine falsche Erklärung rechtliche Folgen haben wird.
\vspace{3cm}

\noindent
\begin{tabular}{p{6cm}p{1cm}p{6cm}}
  Heilbronn, \ThesisDatum & & \\[0.3cm]
  \cline{1-1}\cline{3-3}\\[-0.5cm]
  Ort, Datum              & & \ThesisAutor
\end{tabular}
